\chapter{What is still open}
\label{ch:still-open}

Everything in the previous chapters was fixed, refused, or recorded as a
deliberate decision with a cost. These are the ones that are still live, and they
are listed here rather than in an issue tracker because a defect record that
stops at the fixed ones is a marketing document.

\begin{defect}{The training manifest of the shipped model is marked dirty}
\symptom{The manifest of the training run that produced the currently shipped
model records that the working tree carried uncommitted changes.}
\diagnosis{The model was fitted in the same working session that authored the new
corpus templates. The templates were in the tree and not yet committed at the
moment of fitting, which is the natural order to work in and the wrong order to
record in.}
\resolution{None retroactively. Refitting now to obtain a clean manifest would be
regenerating an artefact to fix a flag, which is the same move the
reproducibility rules forbid for results, and it would produce a different model
from the one every committed evaluation result was measured against.}
\instead{Ignoring it, or refitting. What is done instead is to state it plainly in
three places: the main report's method chapter, its limitations chapter, and
here. What follows from it is bounded and worth being precise about. The model
artefacts are committed and are content addressed into every evaluation run
manifest, so \emph{which} model produced the headline numbers is not in doubt.
The corpus digest in the training manifest matches the digest in every evaluation
manifest, so the model was fitted on the same corpus it was evaluated against.
What is genuinely weaker is the development split evidence, which rests on a run
whose tree state was not pinned to a commit.

The procedural fix is in the future work list and it is one sentence: commit the
corpus change first, then train from the clean tree.}
\end{defect}

\begin{defect}{The evaluation runner still classifies every form twice}
\symptom{Described in full in the language model chapter. The excess disagreement
between a row's predicted label and its finding codes is on the order of two
fields in seven hundred for a non deterministic engine.}
\diagnosis{The runner classifies once directly and once through the audit engine,
and for a non deterministic engine those are separate draws.}
\resolution{Not applied. The fix is to pass the first pass's predictions into the
audit rather than letting the audit reclassify.}
\instead{Fixing it in the phase that measured it. Changing what the runner does
for every engine after seeing the results is a regeneration, and the honest place
for the fix is a phase that re-registers and re-runs.}
\end{defect}

\begin{defect}{A rule in the table that cannot fire}
\symptom{The Japanese postal mark rule, described in the normalisation chapter.}
\diagnosis{Normalisation treats a lone symbol character as a delimiter.}
\resolution{Not applied. Making it reachable means changing what counts as a
token character, which churns every golden snapshot and every committed fixture
expectation.}
\instead{A quiet fix inside an unrelated phase. It is a deliberate change with a
cost and it belongs to a phase that pays the cost on purpose.}
\end{defect}

\begin{defect}{The published package carries no model}
\symptom{An install from the published artefact runs the rule baseline, because
no model bundle travels with it.}
\diagnosis{Model binaries in a package index artefact are a size and licensing
decision that has not been made, and the tool is designed to work without one.}
\resolution{Documented behaviour rather than a silent fallback: the automatic
engine prefers the model when one loads and falls back to the rule baseline with
one printed line when none does.}
\instead{Failing when no model is present, or shipping a model without deciding
the question. The current state has a real cost and it is stated: the default
install is not the default engine that the measurements recommend, and closing
that gap is on the future work list.}
\end{defect}

\begin{defect}{Five cross repository issues, all open}
\symptom{Every issue filed against the evaluation harness is still open. One of
them, publication of that package to an index, is a hard blocker on tagging a
stable release here, because this project's own rule forbids a stable tag while
any dependency is a source reference rather than a released version.}
\diagnosis{The rule is a gate criterion rather than a preference, and it was
written before it became inconvenient.}
\resolution{Not applied. The ledger records each issue with its status and the
workaround, if any, that stands here in the meantime.}
\instead{Vendoring the three functions this project uses and removing the
dependency. That is the cheap move and it would destroy the evidence the boundary
exists to produce, which is the whole reason the dependency is external in the
first place. The rule holds even when holding it is the expensive option, which
is the only circumstance in which a rule tells you anything.}
\end{defect}

\begin{defect}{Nothing here has been measured on a real page}
\symptom{Every accuracy number in this project is a measurement on a corpus this
repository generated.}
\diagnosis{By design, because content from real pages never enters this
repository. The cost is that the strongest objection to every accuracy number is
that a generated corpus has a generator's habits, and the rule table's vocabulary
and the corpus's label text were written by the same person.}
\resolution{Not applied, and it is the single most valuable thing the project
could do next.}
\instead{Relaxing the synthetic data rule. The right shape is a small hand
labelled set of real pages with a published labelling protocol and a published
disagreement rate, which needs no scraping and no stored page content, and it
should arrive with its prediction registered first. The honest registration is
that the rule table's advantage will shrink.}
\end{defect}
